\documentclass[twocolumn,10pt,letterpaper]{article}

% -- Page geometry --
\usepackage[
  top=0.85in,
  bottom=0.75in,
  left=0.65in,
  right=0.65in,
  columnsep=0.25in
]{geometry}

% -- Beautiful typography: Century Schoolbook --
\usepackage[T1]{fontenc}
\usepackage{tgschola}              % TeX Gyre Schola (Century Schoolbook)
\usepackage[scaled=0.82]{beramono} % Bera Mono for code
\usepackage{microtype}             % Microtypographic refinements
\usepackage{setspace}              % Line spacing control
\setstretch{1.08}                  % Slightly open leading for readability

% -- Packages --
\usepackage{xurl}        % Allow URL breaks at any character (load before hyperref)
\usepackage{hyperref}
\usepackage{graphicx}
\usepackage{booktabs}    % Professional tables
\usepackage{tabularx}    % Width-constrained tables
\usepackage{enumitem}    % List customization
\usepackage{fancyhdr}
\usepackage{xcolor}
\usepackage{balance}
\usepackage{stfloats}    % Enable [h] placement for table* in two-column layout
\usepackage{titlesec}    % Section heading customization
\usepackage{listings}    % Code blocks with line wrapping

% -- Paragraph spacing (no indent, modest skip) --
\setlength{\parindent}{0pt}
\setlength{\parskip}{0.4em plus 0.1em minus 0.05em}
\setlength{\emergencystretch}{1em}  % Extra stretch to avoid overfull hboxes in narrow columns

% -- Code listings with line wrapping --
\lstset{
  basicstyle=\small\ttfamily,
  breaklines=true,
  breakatwhitespace=false,
  breakautoindent=true,
  postbreak=\mbox{\textcolor{accentrule}{$\hookrightarrow$}\space},
  columns=fullflexible,
  keepspaces=true,
  backgroundcolor=\color{codebg},
  frame=single,
  rulecolor=\color{codebg},
  framesep=3pt,
  aboveskip=0.5em,
  belowskip=0.5em,
  xleftmargin=4pt,
  xrightmargin=0pt,
}

% -- Colors --
\definecolor{linkblue}{HTML}{1D4ED8}
\definecolor{headergray}{HTML}{1F2937}
\definecolor{rulegray}{HTML}{D1D5DB}
\definecolor{titlebg}{HTML}{111827}
\definecolor{accentrule}{HTML}{6B7280}
\definecolor{brandred}{HTML}{8B2635}
\definecolor{codebg}{HTML}{F3F4F6}    % Very light gray for code backgrounds

% -- Hyperlinks --
\hypersetup{
    colorlinks=true,
    linkcolor=linkblue,
    urlcolor=linkblue,
    citecolor=linkblue,
    pdfstartview=FitH
}

% -- Section numbering: Roman > Alph > Arabic > alph, no decimals --
\renewcommand{\thesection}{\Roman{section}}
\renewcommand{\thesubsection}{\Alph{subsection}}
\renewcommand{\thesubsubsection}{\arabic{subsubsection}}
\renewcommand{\theparagraph}{\alph{paragraph}}
\makeatletter
\@addtoreset{subsection}{section}
\@addtoreset{subsubsection}{subsection}
\@addtoreset{paragraph}{subsubsection}
\makeatother

% -- Section headings --
\titleformat{\section}
  {\large\bfseries\color{titlebg}}
  {\thesection.}{0.5em}{}
  [\vspace{-0.3em}{\color{accentrule}\rule{\columnwidth}{0.5pt}}]
\titleformat{\subsection}
  {\normalsize\bfseries\color{headergray}}
  {\thesubsection.}{0.4em}{}
\titleformat{\subsubsection}
  {\small\bfseries\itshape\color{headergray}}
  {\thesubsubsection.}{0.4em}{}
\titleformat{\paragraph}[runin]
  {\small\bfseries\color{headergray}}
  {\theparagraph)}{0.4em}{}[.\hspace{0.4em}]
\titlespacing*{\section}{0pt}{1.0em}{0.35em}
\titlespacing*{\subsection}{0pt}{0.75em}{0.2em}
\titlespacing*{\subsubsection}{0pt}{0.55em}{0.15em}
\titlespacing*{\paragraph}{0pt}{0.4em}{0em}

% -- Lists: compact, elegant --
\setlist{nosep, leftmargin=1.1em, topsep=0.2em, itemsep=0.05em}
\setlist[itemize]{label={\small\textcolor{accentrule}{\textbullet}}}
\setlist[enumerate]{label={\small\textcolor{headergray}{\arabic*.}}}

% -- Tables: tighter, small text --
\renewcommand{\arraystretch}{1.15}

% -- Header/footer (pages 2+) --
\pagestyle{fancy}
\fancyhf{}
\renewcommand{\headrulewidth}{0.3pt}
\renewcommand{\headrule}{\hbox to\headwidth{\color{rulegray}\leaders\hrule height \headrulewidth\hfill}}
\fancyhead[L]{\footnotesize\color{headergray}\textit{Codex Native Hooks: Verification Before Adopting mickn's Architecture}}
\renewcommand{\footrulewidth}{0.3pt}
\renewcommand{\footrule}{\hbox to\headwidth{\color{rulegray}\leaders\hrule height \footrulewidth\hfill}}
% Brand mark footer: CaseMirror wordmark left, page number right (both burgundy)
\fancyfoot[L]{\footnotesize\textbf{\color{headergray}CaseMirror}}
\fancyfoot[R]{\footnotesize\color{headergray}\thepage}

% -- Page 1 style: footer only, no header rule --
\fancypagestyle{firstpage}{%
  \fancyhf{}%
  \renewcommand{\headrulewidth}{0pt}%
  \renewcommand{\footrulewidth}{0.3pt}%
  \renewcommand{\footrule}{\hbox to\headwidth{\color{rulegray}\leaders\hrule height 0.3pt\hfill}}%
  \fancyfoot[L]{\footnotesize\textbf{\color{headergray}CaseMirror}}%
  \fancyfoot[R]{\footnotesize\color{headergray}\thepage}%
}

% -- Column rule --
\setlength{\columnseprule}{0.15pt}
\renewcommand{\columnseprulecolor}{\color{rulegray}}

% -- Title --
\title{\vspace{-1.8em}\LARGE\bfseries\color{titlebg} Codex Native Hooks: Verification Before Adopting mickn's Architecture\vspace{-0.3em}}
\author{CaseMirror Research \\ \small\itshape\href{https://casemirror.ai}{\textcolor{headergray}{casemirror.ai}}}
\date{{\small\color{headergray} \today}}

\begin{document}

\maketitle
\thispagestyle{firstpage}
\vspace{-0.5em}


\textbf{Generated:} 2026-04-28

\textbf{Topic:} Does the OpenAI Codex CLI actually support \texttt{SessionStart}, \texttt{UserPromptSubmit}, and \texttt{Stop} hooks natively today, or is \texttt{mickn:main}'s rewrite conditional on a future Codex release?


\subsection{Executive Summary}

\textbf{Codex native hooks are real, shipped, and stable.} The OpenAI Codex

CLI exposes \texttt{SessionStart}, \texttt{UserPromptSubmit}, and \texttt{Stop} as first-class

hook events documented at \texttt{developers.openai.com/\allowbreak codex/\allowbreak hooks}. Hooks

were marked stable in \textbf{v0.122.0 (2026-04-20)} via PR \#19012

("Mark codex\_hooks stable"). The locally installed version on this

machine is \textbf{\texttt{codex-\allowbreak cli 0.125.0}} — three releases past the stable

cutoff and well within the supported window.


\texttt{mickn:main}'s v5.0.0 rewrite — which deletes the PTY wrapper, expect

bridge, and queue emitter, replacing them with \texttt{taskmaster-\allowbreak session-\allowbreak start.sh},

\texttt{taskmaster-\allowbreak user-\allowbreak prompt-\allowbreak submit.sh}, and \texttt{taskmaster-\allowbreak stop.sh} — therefore

does \textbf{not} depend on any unreleased Codex feature. It targets shipping,

documented behavior. All three preconditions from the original fork

review (\texttt{docs/\allowbreak upstream-\allowbreak reviews/\allowbreak blader-\allowbreak taskmaster-\allowbreak forks.md} §B1) are

satisfied:


\begin{enumerate}
\setlength{\itemsep}{0pt}
\setlength{\parskip}{0pt}
  \item Codex CLI exposes \texttt{SessionStart}, \texttt{UserPromptSubmit}, and \texttt{Stop} ✅
  \item The native \texttt{Stop} hook supports \texttt{decision: "block"} continuation ✅
  \item \texttt{last\_assistant\_message} is populated on Stop events ✅ (with one
\end{enumerate}
   caveat — see §3)


The answer to the open question is: **proceed with the port. The

wrapper layer is dead weight on Codex 0.122+.** Two caveats worth

gating on are documented in §3 and flow into the punch list.


\subsection{Research Findings}

\subsubsection{Hook events explicitly documented}

The official Codex hooks reference

(\href{https://developers.openai.com/codex/hooks}{developers.openai.com/codex/hooks})

lists six hook events: \texttt{SessionStart}, \texttt{PreToolUse}, \texttt{PermissionRequest},

\texttt{PostToolUse}, \texttt{UserPromptSubmit}, \texttt{Stop}. The three Taskmaster needs

are all present and have dedicated sections.


Hooks are configured via \texttt{\textasciitilde{}/\allowbreak .codex/\allowbreak hooks.json} (user scope) or

\texttt{<repo>/\allowbreak .codex/\allowbreak hooks.json} (project scope), with optional inline

configuration in \texttt{config.toml}. Per-layer hooks are merged, not

overridden — higher-precedence layers add to lower ones. Project-local

hooks only load when the \texttt{.codex/\allowbreak } layer is trusted.


\subsubsection{\texttt{Stop} hook semantics match Claude Code}

The docs explicitly state, for the \texttt{Stop} event:


> "For this event, \texttt{decision: "block"} doesn't reject the turn.

> Instead, it tells Codex to continue and automatically creates a new

> continuation prompt"


The \texttt{reason} field becomes the continuation prompt. This is the same

contract Claude Code's Stop hook uses, which is exactly what

\texttt{taskmaster-\allowbreak stop.sh} needs in order to push a TASKMASTER continuation

prompt back into the same running session.


\texttt{Stop}'s stdin payload includes:


\begin{itemize}
\setlength{\itemsep}{0pt}
\setlength{\parskip}{0pt}
  \item \texttt{turn\_id} — the active Codex turn ID
  \item \texttt{stop\_hook\_active} — whether continuation has already occurred
\end{itemize}
  (the standard guard against infinite re-fire loops; matches Claude's

  field of the same name)

\begin{itemize}
\setlength{\itemsep}{0pt}
\setlength{\parskip}{0pt}
  \item \texttt{last\_assistant\_message} — "Latest assistant message text, if available"
\end{itemize}

The "if available" caveat on \texttt{last\_assistant\_message} is the only

non-trivial parity gap with Claude Code. It is the basis for caveat

(C2) in §3.


\subsubsection{Release timeline}

From the Codex changelog

(\href{https://developers.openai.com/codex/changelog}{developers.openai.com/codex/changelog}):


\vspace{0.4em}
\noindent
\small
\begin{tabularx}{\columnwidth}{@{}XXX@{}}
\toprule
\textbf{Version} & \textbf{Date} & \textbf{Hook-related change} \\
\midrule
v0.116.0 & 2026-03 & Hooks present in experimental form (referenced in issue \#15266 reproductions) \\
v0.122.0 & 2026-04-20 & \texttt{PermissionRequest} hooks added (\#17563); OTEL metrics for hook runs (\#18026) \\
v0.123.0 & 2026-04-23 & Hooks in \texttt{config.toml} / \texttt{requirements.toml} (\#18893); MCP tool support in hooks (\#18385); \textbf{\texttt{codex\_hooks} marked stable (\#19012)} \\
v0.124.0 & 2026-04-23 & \texttt{apply\_patch} emits hooks (\#18391); Bash \texttt{PostToolUse} on \texttt{exec\_command} (\#18888); \textbf{regression: hooks broke at startup if config used map syntax (\#19199)} \\
v0.125.0 & 2026-04 & (locally installed; current) \\
\bottomrule
\end{tabularx}
\vspace{0.4em}

The stable marker landed eight days before this report. mickn's rewrite

(repo timeline aligns with v5.0.0 around the same window) targets the

post-stable surface, not pre-release behavior.


\subsubsection{Known issues that don't block adoption but warrant gating}

**Issue \#15266 — SessionStart + UserPromptSubmit fire simultaneously on

first prompt** (\href{https://github.com/openai/codex/issues/15266}{github.com/openai/codex/issues/15266}).

Filed against v0.116.0 (March 2026). Closed, but the closing

commit/version is not visible in the page content. Behavior described:

on the first prompt of a session, both hooks fire concurrently rather

than \texttt{SessionStart} completing before \texttt{UserPromptSubmit}. On subsequent

prompts, only \texttt{UserPromptSubmit} fires correctly.


Implication for Taskmaster: if \texttt{taskmaster-\allowbreak session-\allowbreak start.sh} writes

state that \texttt{taskmaster-\allowbreak user-\allowbreak prompt-\allowbreak submit.sh} reads (e.g.,

seeding the per-session state file), there's a race on the first

prompt. mickn's \texttt{taskmaster-\allowbreak session-\allowbreak start.sh} is 27 LOC — small enough

to inspect for whether it depends on this ordering. We should verify

on 0.125.0 before merging.


\textbf{Issue \#19199 — v0.124.0 hook config parsing regression}

(\href{https://github.com/openai/codex/issues/19199}{github.com/openai/codex/issues/19199}).

\texttt{codex-\allowbreak cli} failed to start when hooks were configured in

\texttt{config.toml} using map syntax (the documented form). Closed; resolution

version not shown. The local install is 0.125.0, which post-dates the

fix, so this is informational only — but it's a reminder that hook

config schemas are still in flux at the toml-vs-json boundary.


\subsubsection{Third-party confirmation}

Independent projects already shipping against Codex's native hooks:


\begin{itemize}
\setlength{\itemsep}{0pt}
\setlength{\parskip}{0pt}
  \item \textbf{\texttt{hatayama/\allowbreak codex-\allowbreak hooks}} — a hooks runner that reuses Claude
\end{itemize}
  Code's hooks settings against Codex CLI

  (\href{https://github.com/hatayama/codex-hooks}{github.com/hatayama/codex-hooks}).

  Existence of this project confirms the surface is real and

  Claude-Code-compatible enough to be adapted.

\begin{itemize}
\setlength{\itemsep}{0pt}
\setlength{\parskip}{0pt}
  \item \textbf{\texttt{Yeachan-\allowbreak Heo/\allowbreak oh-\allowbreak my-\allowbreak codex} (OmX)} — a Codex enhancement framework
\end{itemize}
  with an active roadmap issue (\#1307) about mapping its hook surfaces

  onto Codex's native hooks, indicating a community migration from

  bespoke wrappers to native is in progress.

\begin{itemize}
\setlength{\itemsep}{0pt}
\setlength{\parskip}{0pt}
  \item \textbf{ArcKit v4} — released March 2026 with first-class Codex hooks
\end{itemize}
  support

  (\href{https://medium.com/arckit/arckit-v4-first-class-codex-and-gemini-support-with-hooks-mcp-servers-and-native-policies-abdf9569e00e}{medium.com/arckit/arckit-v4}).


The pattern across all three: bespoke PTY/wrapper hacks are being

deleted in favor of the native hook surface throughout April 2026.

mickn's rewrite is the same move applied to Taskmaster.


\subsection{Analysis}

The fork review's §B1 set three preconditions for adopting mickn's

wholesale wrapper deletion. All three are satisfied:


\begin{enumerate}
\setlength{\itemsep}{0pt}
\setlength{\parskip}{0pt}
  \item \textbf{Codex CLI exposes SessionStart/UserPromptSubmit/Stop hooks}
\end{enumerate}
   in the version the user is on. Local install: \texttt{codex-\allowbreak cli 0.125.0}.

   Hook events documented since v0.122 stable; we are on 0.125.

   ✅ confirmed.


\begin{enumerate}
\setlength{\itemsep}{0pt}
\setlength{\parskip}{0pt}
  \item \textbf{The native \texttt{Stop} hook supports \texttt{decision: "block"} continuation}
\end{enumerate}
   in the same way Claude Code does. Documented verbatim in

   \texttt{developers.openai.com/\allowbreak codex/\allowbreak hooks}: "doesn't reject the turn.

   Instead, it tells Codex to continue and automatically creates a new

   continuation prompt." ✅ confirmed.


\begin{enumerate}
\setlength{\itemsep}{0pt}
\setlength{\parskip}{0pt}
  \item \textbf{\texttt{last\_assistant\_message} is populated by Codex on stop events.}
\end{enumerate}
   Documented as a Stop-event stdin field, with the qualifier "if

   available." This matches Claude Code's behavior, which also has

   cases where the field is empty (e.g., when the assistant emits no

   message text on stop). ⚠️ confirmed-with-caveat.


The caveat on (3) is meaningful but not blocking. The current fork's

detection is layered: \texttt{last\_assistant\_message} for primary detection,

transcript-grep for explicit \texttt{TASKMASTER\_DONE::<session\_id>} token as

fallback. That layering should survive the port unchanged — the

fallback handles the "if available" gap.


The PTY wrapper, expect bridge, and queue emitter become genuinely

redundant at 0.122+. They cost LOC, complexity, and a \texttt{expect} runtime

dependency. Their only remaining justification was as a portability

floor for Codex versions without native hooks — which now means

versions older than April 2026, a window users will only stay in

deliberately.


The right migration shape mirrors §B1's hedge: keep both code paths,

gate on \texttt{codex -\allowbreak -\allowbreak version} or a feature probe (`codex --help | grep -q

hook\texttt{ or test for the }\textasciitilde{}/.codex/hooks.json` schema), and let

\texttt{install.sh} choose. Default to native on detection, fall back to

wrapper on older Codex. Delete the wrapper path only after a deprecation

window where logs confirm zero installs are using it.


\subsection{Recommendations}

Adopt mickn's native-hooks architecture, but stage it. Two safety

rails make this safe rather than risky:


\begin{enumerate}
\setlength{\itemsep}{0pt}
\setlength{\parskip}{0pt}
  \item \textbf{Version-gated install.} Probe Codex version in \texttt{install.sh}.
\end{enumerate}
   If \texttt{>= 0.122.0}, install native hooks; if \texttt{< 0.122.0} or no codex

   detected, install the existing wrapper path. The user's machine

   (0.125.0) gets the native path automatically.  

\begin{enumerate}
\setlength{\itemsep}{0pt}
\setlength{\parskip}{0pt}
  \item \textbf{Keep the wrapper path on disk.} Don't \texttt{git rm} the PTY wrapper,
\end{enumerate}
   expect bridge, or their tests in the same PR. Mark them

   \texttt{legacy/\allowbreak }-prefixed and have \texttt{install.sh} install from \texttt{legacy/\allowbreak }

   when version-gated. Plan to remove them after one minor release if

   no one reports using them.


The \texttt{last\_assistant\_message} + transcript-token layered detection

already in the fork is the right pattern and ports cleanly. Do not

collapse to a single detection mode.


\subsection{Punch List (for \texttt{/\allowbreak mei})}

Each item is a self-contained adoption decision. Numbered for

priority. Phase tags only — no time estimates.


\begin{enumerate}
\setlength{\itemsep}{0pt}
\setlength{\parskip}{0pt}
  \item \textbf{[Phase 1, HIGH] Add Codex version probe to \texttt{install.sh}.}
\end{enumerate}
   Detect \texttt{codex -\allowbreak -\allowbreak version} and parse semver; expose as

   \texttt{$CODEX\_HOOKS\_NATIVE} (true if \texttt{>= 0.122.0}). Touches only

   \texttt{install.sh}. No behavior change yet — just the detection.


\begin{enumerate}
\setlength{\itemsep}{0pt}
\setlength{\parskip}{0pt}
  \item **[Phase 1, HIGH] Port \texttt{taskmaster-\allowbreak session-\allowbreak start.sh} from
\end{enumerate}
   \texttt{mickn:main}.** 27 LOC. Place at \texttt{hooks/\allowbreak taskmaster-\allowbreak session-\allowbreak start.sh}.

   Verify it does not depend on completing-before-\texttt{UserPromptSubmit}

   ordering (issue \#15266). If it does, add a state-file lock that

   both hooks honor.


\begin{enumerate}
\setlength{\itemsep}{0pt}
\setlength{\parskip}{0pt}
  \item \textbf{[Phase 1, HIGH] Port \texttt{taskmaster-\allowbreak stop.sh} from \texttt{mickn:main}.}
\end{enumerate}
   356 LOC. Replaces the wrapper's stop-detection role. Must emit

   \texttt{decision: "block"} JSON with the shared compliance prompt as

   \texttt{reason}. Reuses \texttt{taskmaster-\allowbreak compliance-\allowbreak prompt.sh}. Keep the

   \texttt{last\_assistant\_message} → transcript-token fallback layered

   detection.


\begin{enumerate}
\setlength{\itemsep}{0pt}
\setlength{\parskip}{0pt}
  \item **[Phase 1, HIGH] Port \texttt{taskmaster-\allowbreak user-\allowbreak prompt-\allowbreak submit.sh} from
\end{enumerate}
   \texttt{mickn:main}.** 107 LOC. Implements per-turn external user prompt

   capture with \texttt{<hook\_prompt>} filtering. This is pattern A1 from the

   fork review and the highest-value behavioral upgrade.


\begin{enumerate}
\setlength{\itemsep}{0pt}
\setlength{\parskip}{0pt}
  \item \textbf{[Phase 1, MEDIUM] Move existing wrapper artifacts to \texttt{legacy/\allowbreak }.}
\end{enumerate}
   \texttt{hooks/\allowbreak inject-\allowbreak continue-\allowbreak codex.sh}, \texttt{hooks/\allowbreak run-\allowbreak codex-\allowbreak expect-\allowbreak bridge.exp},

   \texttt{run-\allowbreak taskmaster-\allowbreak codex.sh}, plus their tests in

   \texttt{tests/\allowbreak inject-\allowbreak continue-\allowbreak codex.test.sh} etc. Update \texttt{install.sh} to

   choose \texttt{hooks/\allowbreak } (native) or \texttt{legacy/\allowbreak } (wrapper) based on the

   version probe.


\begin{enumerate}
\setlength{\itemsep}{0pt}
\setlength{\parskip}{0pt}
  \item \textbf{[Phase 1, MEDIUM] Write \texttt{\textasciitilde{}/\allowbreak .codex/\allowbreak hooks.json} template} in
\end{enumerate}
   \texttt{install.sh} mapping the three event names to the three new

   \texttt{hooks/\allowbreak taskmaster-\allowbreak *.sh} scripts. Merge-safe with any existing

   user hooks (don't clobber).


\begin{enumerate}
\setlength{\itemsep}{0pt}
\setlength{\parskip}{0pt}
  \item \textbf{[Phase 2, MEDIUM] Add a feature smoke test:} create `taskmaster
\end{enumerate}
   selftest --codex-hooks` that fires a no-op session, asserts each of

   the three hooks executed via state-file markers, and reports OK/FAIL.

   Exercised in \texttt{install.sh -\allowbreak -\allowbreak verify} and CI.


\begin{enumerate}
\setlength{\itemsep}{0pt}
\setlength{\parskip}{0pt}
  \item **[Phase 2, LOW] Document the version gate in
\end{enumerate}
   \texttt{docs/\allowbreak SPEC.md}.** Single section explaining the dual code paths,

   the 0.122.0 cutover, the issue-\#15266 race awareness, and the

   deprecation plan for \texttt{legacy/\allowbreak }.


\begin{enumerate}
\setlength{\itemsep}{0pt}
\setlength{\parskip}{0pt}
  \item \textbf{[Phase 3, LOW] Sunset \texttt{legacy/\allowbreak } after one minor release} if no
\end{enumerate}
   reports of installs using it. Track via an \texttt{install.sh}

   instrumentation line that logs which path was selected. Drop after

   evidence justifies it.


\begin{enumerate}
\setlength{\itemsep}{0pt}
\setlength{\parskip}{0pt}
  \item \textbf{[Phase 3, LOW] Watch upstream issue \#15266} for a definitive
\end{enumerate}
    fix-version. If the simultaneous-fire race is confirmed fixed in

    a known version, tighten \texttt{$CODEX\_HOOKS\_NATIVE} lower bound to

    that version and remove any race-mitigation lock added in (2).


\subsection{Sources}

\begin{itemize}
\setlength{\itemsep}{0pt}
\setlength{\parskip}{0pt}
  \item \href{https://developers.openai.com/codex/hooks}{Hooks – Codex | OpenAI Developers} — primary reference; documents all six hook events including \texttt{SessionStart}, \texttt{UserPromptSubmit}, \texttt{Stop}. Contains the verbatim specification of \texttt{decision: "block"} for \texttt{Stop} and the \texttt{last\_assistant\_message} field.
  \item \href{https://developers.openai.com/codex/changelog}{Changelog – Codex | OpenAI Developers} — release timeline confirming hooks went stable in v0.122.0 (April 20, 2026) via PR \#19012 "Mark codex\_hooks stable."
  \item \href{https://github.com/openai/codex/issues/15266}{Issue \#15266 — UserPromptSubmit and SessionStart hooks fire simultaneously on first prompt} — known caveat, filed v0.116.0 (March 2026), closed.
  \item \href{https://github.com/openai/codex/issues/19199}{Issue \#19199 — codex-cli 0.124.0 fails to start when hook config is present and codex\_hooks is enabled} — informational; not a blocker on 0.125.0.
  \item \href{https://github.com/openai/codex/pull/14867}{PR \#14867 — hooks: use a user message > developer message for prompt continuation} — early-development context for hook continuation semantics.
  \item \href{https://github.com/openai/codex/pull/15118}{PR \#15118 — hooks: turn\_id extension for Stop \& UserPromptSubmit} — confirms \texttt{turn\_id} field added to the stdin payload of these specific hooks.
  \item \href{https://github.com/openai/codex/discussions/2150}{Discussion \#2150 — Hook would be a great feature} — historical context for how hooks landed.
  \item \href{https://github.com/hatayama/codex-hooks}{hatayama/codex-hooks} — third-party hooks runner reusing Claude Code's hooks settings against Codex; corroborates Claude-compatible surface.
  \item \href{https://github.com/Yeachan-Heo/oh-my-codex/issues/1307}{Yeachan-Heo/oh-my-codex issue \#1307 — roadmap: map OMC hook surfaces onto OMX native Codex hooks} — community-wide migration pattern from bespoke wrappers to native hooks.
  \item \href{https://medium.com/arckit/arckit-v4-first-class-codex-and-gemini-support-with-hooks-mcp-servers-and-native-policies-abdf9569e00e}{ArcKit v4: First-Class Codex and Gemini Support with Hooks} — independent third-party adoption of the same hook surface in March 2026.
  \item Local fork review: \texttt{docs/\allowbreak upstream-\allowbreak reviews/\allowbreak blader-\allowbreak taskmaster-\allowbreak forks.md} §B1 — the three preconditions verified by this report.
  \item Local Codex install: \texttt{codex-\allowbreak cli 0.125.0} (\texttt{/\allowbreak usr/\allowbreak local/\allowbreak bin/\allowbreak codex}) — three releases past the stable cutoff.
\end{itemize}


\balance
\end{document}
